% Data-derived report section; editorial comments are preserved in this template.
\section{Measurement Definitions}
\subsection{Source Reference and Conversion Gain}
All results use the same available input power reference and differential IF
output. For a source resistance $R_s$, an open-circuit sinusoidal peak voltage
$V_{\mathrm{open,pk}}$ gives
\begin{equation}
P_{\mathrm{in,av}}=\frac{V_{\mathrm{open,pk}}^2}{8R_s}.
\label{eq:pin}\end{equation}
At the small-signal test power of $-55$~dBm, the open-circuit peak is
1.12468~mV. The analogLib port uses matched-load voltage entries: enter half
this value (0.56234~mV), with $R_s=50~\Omega$ inside the port. The active
port's DC entry is 0.325~V to obtain 0.65~V open circuit; the passive entry is
zero. Its PAC magnitude is 0.5~V for a 1~V open-circuit reference. Do not add a
second source resistor. LO legs use 0.2~V sine peaks and phases of
$0^\circ$ and $180^\circ$.

For differential IF peak amplitude $V_{\mathrm{IF,pk}}$ and one physical
branch resistor $R_L$, the reported gain and load power are
\begin{align}
G_V &=20\log_{10}\!\left(\frac{V_{\mathrm{IF,pk}}}{V_{\mathrm{open,pk}}}\right),
\label{eq:gv}\\
P_{\mathrm{IF,load}}&=\frac{V_{\mathrm{IF,pk}}^2}{4R_L},
\label{eq:pout}\\
G_T&=10\log_{10}\!\left(\frac{P_{\mathrm{IF,load}}}{P_{\mathrm{in,av}}}\right).
\label{eq:gt}\end{align}
Voltage conversion loss is $-G_V$ and transducer power loss is $-G_T$.
The different source/load impedances mean that voltage gain and power gain
must not be interchanged.

\subsection{Transient and Linearity Measurements}
Single-tone measurements use 5~ns settling followed by a 20~ns observation
window at a 5~ps strobe interval (4000 samples). A 20~ns settling plus 100~ns
window checks steady-state convergence. Spectra use a rectangular DFT on the
settled differential output. The time-domain plots retain the full waveform,
including LO/RF feedthrough, rather than an extracted IF sinusoid.

The input 1~dB compression point is the first crossing of the low-power gain
minus 1~dB in a fourteen-point input-power campaign. The reported midpoint
has a half-bracket resolution of 0.078125~dB. This is sweep resolution, not a
statistical confidence interval.

Equal tones at 5.79 and 5.81~GHz produce IF fundamentals at 190 and 210~MHz,
and third-order products at 170 and 230~MHz. Two-tone measurements use a
20~ns settling time and 100~ns window (20000 samples; 10~MHz bins).
For each resolved product,
\begin{equation}
P_{\mathrm{IIP3}}=P_{\mathrm{in,tone}}
+10\log_{10}\!\left(\frac{A_{\mathrm{IF}}}{A_{\mathrm{IM3}}}\right),
\label{eq:iip3}\end{equation}
where the powers are in dBm and the amplitudes use the same peak-voltage
reference. The accepted interval is $-40$, $-35$ and $-30$~dBm per tone;
fundamental and IM3 slopes must approach one and three, respectively.
The lower of the two product intercepts is reported. A repeat at $-35$~dBm
uses effective reltol $10^{-7}$ and a 1~ps maximum step, versus $10^{-6}$ and
2~ps for the main campaign. Empty-bin margins test numerical resolution,
not a stochastic noise floor.

\subsection{Periodic Noise}
Shooting periodic steady-state analysis uses the 5.6~GHz LO, 63 harmonics,
and 5~ns stabilization. RF transient amplitudes are zero; RF DC bias, source
resistance noise and PAC excitation remain enabled. Source noise temperature
is 290~K and circuit temperature is $27^\circ$C. Periodic noise is evaluated
at 100, 200 and 300~MHz while checking 15, 31 and 63 folded sidebands.

Let $S_{v,\mathrm{out}}$ be total differential output voltage-noise PSD and
$G_{\mathrm{USB}}$, $G_{\mathrm{image}}$ the desired and image voltage gains
referenced to the open-circuit source. The noise factors are
\begin{align}
F_{\mathrm{SSB}}&=\frac{S_{v,\mathrm{out}}}{4kT_0R_s|G_{\mathrm{USB}}|^2},
\label{eq:ssb}\\
F_{\mathrm{DSB}}&=\frac{S_{v,\mathrm{out}}}
{4kT_0R_s\bigl(|G_{\mathrm{USB}}|^2+|G_{\mathrm{image}}|^2\bigr)}.
\label{eq:dsb}\end{align}
Noise figure is $10\log_{10}F$. At 200~MHz IF, the desired and image RF
frequencies are 5.8 and 5.4~GHz. The DSB result uses both measured gains;
it is not defined by subtracting 3~dB from SSB NF. Both branch-resistor noise
contributions are included. Ideal LO-driver noise and power are excluded.
